\documentclass[11pt,a4paper]{article}

% =========================================================
% COMMON PACKAGES
% =========================================================
\usepackage{lmodern}
\usepackage[T1]{fontenc}
\usepackage[utf8]{inputenc}
\usepackage[english]{babel}

\usepackage{amsmath,amssymb,amsfonts,amsthm,mathtools}
\usepackage{bm}
\usepackage{microtype}
\usepackage{enumitem}
\usepackage{csquotes}
\usepackage{mathrsfs}
\usepackage{thmtools}
\usepackage{hyperref}
\usepackage[nameinlink,capitalize]{cleveref}
\usepackage{geometry}

\geometry{
	a4paper,
	margin=2.8cm
}

% =========================================================
% THEOREM STYLES
% =========================================================
\numberwithin{equation}{section}

\declaretheoremstyle[
headfont=\bfseries,
bodyfont=\itshape,
spaceabove=8pt,
spacebelow=8pt
]{plainstyle}

\declaretheoremstyle[
headfont=\bfseries,
bodyfont=\normalfont,
spaceabove=8pt,
spacebelow=8pt
]{defstyle}

\declaretheoremstyle[
headfont=\itshape,
bodyfont=\normalfont,
spaceabove=4pt,
spacebelow=4pt
]{remarkstyle}

\declaretheorem[style=plainstyle,name=Theorem]{theorem}
\declaretheorem[style=plainstyle,name=Proposition,sibling=theorem]{proposition}
\declaretheorem[style=plainstyle,name=Lemma,sibling=theorem]{lemma}
\declaretheorem[style=plainstyle,name=Corollary,sibling=theorem]{corollary}

\declaretheorem[style=defstyle,name=Definition,sibling=theorem]{definition}
\declaretheorem[style=defstyle,name=Example,sibling=theorem]{example}

\declaretheorem[style=remarkstyle,name=Remark,sibling=theorem]{remark}

% Cleveref names
\crefname{theorem}{Theorem}{Theorems}
\crefname{proposition}{Proposition}{Propositions}
\crefname{lemma}{Lemma}{Lemmas}
\crefname{corollary}{Corollary}{Corollaries}
\crefname{definition}{Definition}{Definitions}
\crefname{example}{Example}{Examples}
\crefname{remark}{Remark}{Remarks}
\crefname{equation}{Equation}{Equations}
\crefname{section}{Section}{Sections}

% =========================================================
% COMMON MACROS FOR THE TWIN-SPACES SERIES
% =========================================================

% Sets and fields
\newcommand{\R}{\mathbb{R}}
\newcommand{\C}{\mathbb{C}}
\newcommand{\Z}{\mathbb{Z}}
\newcommand{\N}{\mathbb{N}}

% Operators
\DeclareMathOperator{\End}{End}
\DeclareMathOperator{\Hom}{Hom}
\DeclareMathOperator{\id}{id}
\DeclareMathOperator{\Ad}{Ad}
\DeclareMathOperator{\Aut}{Aut}
\DeclareMathOperator{\Ker}{Ker}
\DeclareMathOperator{\Img}{Im}
\DeclareMathOperator{\HH}{HH}

% Twin-space notation
\newcommand{\TwinSpace}{(E,\alpha,G)}
\newcommand{\Twin}{\mathcal{T}}
\newcommand{\TwinOrbit}{O}
\newcommand{\LinearSubgroup}{L}

% Cohomology notation
\newcommand{\Cochain}{C}
\newcommand{\diff}{\delta}

% Misc
\newcommand{\abs}[1]{\left\lvert #1 \right\rvert}
\newcommand{\norm}[1]{\left\lVert #1 \right\rVert}
\newcommand{\bracket}[1]{\left\langle #1 \right\rangle}

% =========================================================
% METADATA FOR ARTICLE A
% =========================================================

\title{\textbf{Deformation Cohomology of Twin Operations}}
\author{Jocelyn Nembe}
\date{\today}

% =========================================================
% DOCUMENT
% =========================================================

\begin{document}
	
	\maketitle
	
	\begin{abstract}
		We develop a complete deformation cohomology theory for \emph{twin operations},
		a group--action--based formalism introduced in Article~I.
		Given a set $E$ equipped with a binary operation $\alpha : E\times E \to E$
		and a group $G$ acting on $E$, we consider the full family of 
		\emph{twin operations}
		\[
		\alpha_g(x,y) = g^{-1}\cdot \alpha(g\cdot x,\ g\cdot y), \qquad g\in G,
		\]
		which forms an orbit of operations canonically identified with the homogeneous space $\LinearSubgroup\backslash G$ of \emph{right} cosets,
		where $\LinearSubgroup=\{g\in G : \alpha_g=\alpha\}$ is the linear subgroup.
		This orbit perspective---viewing a symmetry as generating a \emph{space} of operations rather than a single invariant one---is the key structural ingredient that makes deformation cohomology possible in this setting.
		
		We show that this twin--space formalism admits a Hochschild-type cohomology.
		In particular:
		(i) $\HH^1(G,\End(E))$ classifies infinitesimal deformations of the twin system $\TwinSpace$;
		(ii) $\HH^2(G,\End(E))$ contains the obstructions to extending first-order deformations;
		(iii) rigidity of the twin system under $G$-equivariant deformations is equivalent to the vanishing of $\HH^1$;
		(iv) explicit computations are provided for cyclic, abelian, dihedral, and compact Lie groups, as well as matrix algebras;
		(v) a Hochschild--Serre spectral sequence is established.
		
		This yields a deformation--obstruction theory for twin spaces that parallels, but is conceptually distinct from, classical Gerstenhaber deformation theory.
	\end{abstract}
	
	\tableofcontents
	
	%===========================================================
	\section{Introduction}
	\label{sec:intro}
	
	Classical deformation theory investigates how algebraic structures vary when perturbed,
	and was developed in foundational works of Gerstenhaber~\cite{Gerstenhaber1964},
	Nijenhuis--Richardson~\cite{NijenhuisRichardson1967},
	Drinfeld~\cite{Drinfeld1986}, and Kontsevich~\cite{Kontsevich2003}.
	A central theme is that infinitesimal deformations are controlled by cohomology groups. The prototype is the deformation theory of an associative algebra $(A,\mu)$: a first-order perturbation $\mu\mapsto\mu+t\,\mu_1$ remains associative to order $t$ exactly when $\mu_1$ is a Hochschild $2$-cocycle, two such perturbations are equivalent exactly when they differ by a $2$-coboundary, and the obstruction to integrating $\mu_1$ to second order lives in $HH^3(A,A)$. Thus $HH^2$ classifies and $HH^3$ obstructs. The theory developed below follows the same script---cocycles classify, coboundaries trivialise, the next degree obstructs---but the object being deformed is not a single multiplication: it is an entire \emph{orbit of operations} produced by a symmetry group, and the controlling cohomology is correspondingly different.
	
	In Article~I we introduced \emph{twin operations}, a new algebraic construction attached to a group action.
	Let $E$ be a set with a left action of a group $G$, and let $\alpha : E\times E \to E$ be a binary operation.
	Rather than focus only on $G$-invariant operations, we consider the entire orbit of $\alpha$ under the induced action
	\begin{equation}\label{eq:twin-def}
		\alpha_g(x,y) := g^{-1}\cdot \alpha(g\cdot x,\ g\cdot y), \qquad g\in G.
	\end{equation}
	The resulting family of operations
	\[
	\TwinOrbit = \{\alpha_g : g\in G\}
	\]
	is not arbitrary: there is a canonical bijection
	\[
	\TwinOrbit \;\cong\; \LinearSubgroup\backslash G \qquad(\text{right cosets}),\qquad \abs{\TwinOrbit}=[G:\LinearSubgroup],
	\]
	where 
	\[
	\LinearSubgroup := \{ g\in G : \alpha_g=\alpha \}
	\]
	is the \emph{linear subgroup}.  
	We call this orbit of operations, equipped with the underlying action, a \emph{twin space}.

	Before developing any cohomology it is worth seeing the orbit $\TwinOrbit$ explicitly in a case where the operations genuinely move.

	\begin{example}[The translated additions]\label{ex:translated}
		Let $E=\R$, let $\alpha(x,y)=x+y$ be ordinary addition, and let $G=\mathrm{Aff}(\R)=\{g_{a,b}:x\mapsto ax+b\mid a\neq 0\}$
		act on $\R$. Using $g_{a,b}^{-1}(u)=(u-b)/a$ we compute
		\[
			x\,\alpha_{g_{a,b}}\,y \;=\; g_{a,b}^{-1}\!\big((ax+b)+(ay+b)\big)
			\;=\; \frac{a(x+y)+2b-b}{a} \;=\; x+y+\tfrac{b}{a}.
		\]
		So $\alpha_{g_{a,b}}$ is ``addition shifted by $c=b/a$'', and the twin orbit is the one-parameter
		family $\TwinOrbit=\{(x,y)\mapsto x+y+c : c\in\R\}$. The linear subgroup is
		$\LinearSubgroup=\{g_{a,0}\}$, the pure scalings (for which $c=0$), and indeed
		$\TwinOrbit\cong\LinearSubgroup\backslash\mathrm{Aff}(\R)\cong\R$ through $g_{a,b}\mapsto b/a$.
		Each $\alpha_{g}$ is isomorphic to $\alpha$ (it is carried onto $\alpha$ by $g$), yet the operations
		are pairwise distinct on the fixed set $\R$: a single operation has been spread by the symmetry into
		a whole line of operations. Deformation theory asks how such a family can be perturbed
		$G$-equivariantly, and which perturbations are genuinely new.
	\end{example}
	
	This viewpoint---from a single operation to a \emph{space of operations generated by symmetry}---has three structural consequences that are central to this article:
	
	\begin{enumerate}[label=(\roman*)]
		\item \textbf{New deformation complex.}  
		The deformation problem concerns the entire twin space $\TwinOrbit\cong \LinearSubgroup\backslash G$ rather than $\alpha$ alone.
		This leads naturally to a Hochschild-type cochain complex 
		\[
		\Cochain^\bullet(G,\End(E))
		\]
		governing deformations of the twin structure $\TwinSpace$.
		
		\item \textbf{Geometric classification mechanism.}  
		Because $\TwinOrbit$ is a homogeneous space $\LinearSubgroup\backslash G$, classifying compatible operations becomes
		a geometric and group--theoretic problem of describing subgroups of $G$ containing $\LinearSubgroup$.
		
		\item \textbf{Deformation--transport duality.}  
		Deforming $\alpha$ means deforming the \emph{transport rule}
		\[
		\alpha \mapsto \alpha_g,
		\]
		not the operation alone.  
		Infinitesimal deformations correspond to infinitesimal variations in this transport,
		and are measured precisely by $\HH^1(G,\End(E))$.
	\end{enumerate}
	
	The purpose of this article is to develop a full deformation cohomology theory for twin spaces:
	we construct a Hochschild-type complex, prove that $\HH^1$ classifies infinitesimal deformations,
	identify obstruction classes in $\HH^2$, provide cohomological rigidity criteria,
	and compute these groups in a range of examples, including a Hochschild--Serre spectral sequence
	of Hochschild--Serre type~\cite{HochschildSerre1953}.
	
	%===========================================================
	\section{Twin Operations and Group Actions}
	\label{sec:twin}
	
	Let $G$ act on $E$ from the left.  
	Given a binary operation $\alpha : E\times E \to E$, we define
	\begin{equation}\label{eq:twin-op}
		\alpha_g(x,y) := g^{-1}\cdot \alpha(g\cdot x,\ g\cdot y).
	\end{equation}
	
	\begin{definition}[Linear subgroup]
		The \emph{linear subgroup} of $\TwinSpace$ is
		\[
		\LinearSubgroup := \{g\in G : \alpha_g = \alpha\}.
		\]
	\end{definition}
	
	\begin{proposition}\label{prop:orbit}
		The map
		\[
		\LinearSubgroup\backslash G \longrightarrow \TwinOrbit, \qquad \LinearSubgroup g \mapsto \alpha_g,
		\]
		is a well-defined bijection.
		In particular, $\TwinOrbit$ is canonically isomorphic to the homogeneous space $\LinearSubgroup\backslash G$ of \emph{right} cosets, and $\abs{\TwinOrbit}=[G:\LinearSubgroup]$.
	\end{proposition}
	
	\begin{proof}
		Surjectivity is immediate from the definition $\TwinOrbit=\{\alpha_g : g\in G\}$.
		For injectivity, suppose $\alpha_g=\alpha_h$.
		Then for all $x,y\in E$,
		\[
		g^{-1}\cdot \alpha(g\cdot x,g\cdot y)
		=
		h^{-1}\cdot \alpha(h\cdot x,h\cdot y).
		\]
		Applying $g$ to both sides gives
		\[
		\alpha(g\cdot x,g\cdot y)
		=
		gh^{-1}\cdot \alpha(h\cdot x,h\cdot y).
		\]
		Substituting $(x,y)\mapsto(h^{-1}\cdot x,h^{-1}\cdot y)$ yields
		\[
		\alpha(gh^{-1}\cdot x,gh^{-1}\cdot y)
		=
		gh^{-1}\cdot \alpha(x,y),
		\]
		so $gh^{-1}\in \LinearSubgroup$, that is, $\LinearSubgroup g=\LinearSubgroup h$.
		Conversely, if $\LinearSubgroup g=\LinearSubgroup h$ then $gh^{-1}\in\LinearSubgroup$ and reversing the steps gives $\alpha_g=\alpha_h$; hence $\LinearSubgroup g\mapsto\alpha_g$ is well defined and injective.
	\end{proof}
	
	\begin{remark}[Right cosets and group structure]\label{rem:right-cosets}
		The identification is with the space of \emph{right} cosets $\LinearSubgroup\backslash G$,
		not $G/\LinearSubgroup$: the assignment $g\mapsto\alpha_g$ is a \emph{right} action, since a
		direct computation gives $\alpha_{hg}=(\alpha_h)_g$ --- apply the defining formula twice, $(\alpha_h)_g(x,y)=g^{-1}h^{-1}\,\alpha(hg\cdot x,\,hg\cdot y)=\alpha_{hg}(x,y)$ --- so the stabilizer of $\alpha$ is
		$\LinearSubgroup$ and the orbit is $\LinearSubgroup\backslash G$ by orbit--stabilizer.
		Consequently $\abs{\TwinOrbit}=[G:\LinearSubgroup]$ unconditionally, and $\TwinOrbit$ is a
		homogeneous $G$-set. It carries a compatible group structure $\TwinOrbit\cong G/\LinearSubgroup$
		precisely when $\LinearSubgroup\trianglelefteq G$ (e.g.\ whenever $G$ is abelian); in general
		$\LinearSubgroup$ need not be normal. All cohomological results below depend only on the
		$G$-module structure of $M=\End(E)$ and are unaffected by this distinction.
	\end{remark}

	This orbit description is the geometric backbone of twin spaces.
	
	%===========================================================
	\section{The Hochschild-Type Complex for Twin Operations}
	\label{sec:hoch}
	
	Set $M:=\End(E)$ and let $G$ act on $M$ by
	\begin{equation}\label{eq:End-action}
		(g\cdot A)(x) := g\cdot A(g^{-1}\cdot x), \qquad A\in M,\ g\in G.
	\end{equation}
	
	The action~\eqref{eq:End-action} is the natural \emph{conjugation} action of $G$ on linear
	endomorphisms: $g$ transports a state back by $g^{-1}$, applies $A$, and transports the result forward
	by $g$. It makes $M=\End(E)$ a left $G$-module, and the complex defined below is precisely the
	inhomogeneous bar complex computing the group cohomology $H^\bullet(G,M)$. The deformation-theoretic
	dictionary is then transparent: a $1$-cochain $f:G\to M$ records, for each symmetry $g$, the
	infinitesimal endomorphism by which the transported operation $\alpha_g$ is allowed to drift; the
	cocycle equation $\diff f=0$ says these drifts are compatible with composition of symmetries; and the
	coboundaries are exactly the drifts induced by an infinitesimal change of coordinates on $E$. This is
	the ``deformation--transport duality'' announced in \cref{sec:intro}, and we make it precise in
	\cref{sec:low-degree}.

	\begin{definition}[Cochains]
		For $n\ge 0$, the group of $n$-cochains is
		\[
		\Cochain^n(G,M) := \{ f : G^n \to M \}.
		\]
	\end{definition}
	
	\begin{definition}[Hochschild-type differential]
		The differential $\diff : \Cochain^n(G,M)\to \Cochain^{n+1}(G,M)$ is given by
		\begin{align*}
			(\diff f)(g_1,\dots,g_{n+1}) 
			&= g_1\cdot f(g_2,\dots,g_{n+1})
			\\&\quad + \sum_{i=1}^n (-1)^i f(g_1,\dots,g_ig_{i+1},\dots,g_{n+1})
			\\&\quad + (-1)^{n+1} f(g_1,\dots,g_n).
		\end{align*}
	\end{definition}
	
	\begin{proposition}
		$\diff^2=0$.
	\end{proposition}
	
	\begin{proof}
		Write $\diff=\sum_{i=0}^{n+1}(-1)^i d_i$, where $d_0 f=g_1\cdot f(g_2,\dots,g_{n+1})$, the map $d_i$ ($1\le i\le n$) contracts the adjacent pair $g_ig_{i+1}$, and $d_{n+1}$ drops the last argument. These satisfy the simplicial identities $d_i d_j=d_{j-1}d_i$ for $i<j$: for the contraction maps they hold by associativity of the group law, and the two identities involving $d_0$ use the module axiom $g_1\cdot(g_2\cdot A)=(g_1g_2)\cdot A$. Substituting into $\diff^2=\sum_{i,j}(-1)^{i+j}d_id_j$ and re-indexing, the terms cancel in pairs.
		Each term of the form $g_1g_2\cdot f(\cdots)$ appears twice with opposite sign,
		and similarly for all intermediate combinations.
		One recovers exactly the usual cancellation in the bar complex for group cohomology.
	\end{proof}
	
	\begin{definition}[Twin deformation cohomology]
		We define
		\[
		\HH^n(G,M) := \frac{\Ker(\diff : \Cochain^n\to \Cochain^{n+1})}{\Img(\diff : \Cochain^{n-1}\to \Cochain^n)}.
		\]
		This is the Hochschild-type cohomology governing deformations of the twin system $\TwinSpace$.
	\end{definition}
	
	%===========================================================
	\section{Low-Degree Cohomology}
	\label{sec:low-degree}
	
	We describe $\HH^0$, $\HH^1$, and $\HH^2$ and their deformation-theoretic meaning.
	
	\subsection{$\HH^0$: invariant endomorphisms}
	
	A $0$-cochain is an element $A\in M$. Then
	\[
	(\diff A)(g) = g\cdot A - A.
	\]
	Thus
	\[
	\HH^0(G,M) = M^G := \{A\in M : g\cdot A = A\ \forall g\in G\},
	\]
	the subspace of endomorphisms commuting with the $G$-action.
	
	\subsection{$\HH^1$: infinitesimal deformations}
	
	We consider a formal deformation
	\[
	\alpha_t(x,y) = \alpha(x,y) + t\,\psi_1(x,y) + O(t^2).
	\]
	
	\begin{definition}[First-order $G$-equivariance]
		We say that $\alpha_t$ is $G$-equivariant to first order if
		\[
		g\cdot \alpha_t(x,y) - \alpha_t(g\cdot x,g\cdot y) = O(t^2)
		\]
		for all $g\in G$.
	\end{definition}
	
	Expanding in $t$, and using \eqref{eq:twin-op}, we obtain
	\[
	g\cdot\psi_1(x,y) - \psi_1(g\cdot x,g\cdot y) = 0,
	\]
	which is precisely the cocycle condition $\diff f=0$ for the $1$-cochain $f\in \Cochain^1(G,M)$ that
	records the infinitesimal variation of the transport rule. (It is essential to keep the two roles
	distinct: $\psi_1$ perturbs the \emph{operation} as a bilinear map, while $f$ is its shadow on the
	\emph{transport} $g\mapsto\alpha_g$, valued in $\End(E)$. It is the transport---equivalently the
	equivariance of the family $\{\alpha_g\}$---whose deformations are governed by $H^\bullet(G,\End(E))$,
	which is why the cochains take values in $\End(E)$ rather than in bilinear maps.)
	
	\begin{theorem}\label{thm:HH1-deformations}
		Infinitesimal $G$-equivariant deformations of the twin system $\TwinSpace$ are classified,
		up to equivalence, by $\HH^1(G,M)$.
	\end{theorem}
	
	\begin{proof}
		Let $f\in \Cochain^1(G,M)$ encode the infinitesimal change in the transport rule
		$\alpha\mapsto \alpha_g$. The first-order $G$-equivariance condition is exactly $\diff f=0$.
		Two deformations differing by a change of variables
		$x\mapsto x + t\,\varphi(x)$ correspond to $f$ and $f+\diff\varphi$, so they define the same
		class in $\HH^1$. Thus equivalence classes of infinitesimal deformations are in bijection with
		$\HH^1(G,M)$.
	\end{proof}
	
	\subsection{$\HH^2$: obstructions}
	
	To extend a first-order deformation to second order
	\[
	\alpha_t=\alpha + t\psi_1 + t^2\psi_2 + O(t^3)
	\]
	one obtains an equation of the form
	\[
	\diff\psi_2 = \Theta(\psi_1)
	\]
	for a certain 2-cochain $\Theta(\psi_1)$, determined quadratically from $\psi_1$.
	The consistency condition is that $\Theta(\psi_1)$ is a coboundary.
	
	\begin{theorem}\label{thm:HH2-obstruction}
		The obstruction to extending an infinitesimal deformation to second order
		is a cohomology class in $\HH^2(G,M)$. The deformation extends
		if and only if this class vanishes.
	\end{theorem}
	
	\begin{proof}
		The second-order $G$-equivariance condition has the form
		\[
		g\cdot \alpha_t(x,y) - \alpha_t(g\cdot x,g\cdot y) = O(t^3),
		\]
		which at order $t^2$ yields an equation
		\[
		g\cdot \psi_2(x,y) - \psi_2(g\cdot x,g\cdot y) = \Xi(\psi_1)(g;x,y),
		\]
		where the right-hand side is an explicit quadratic expression in $\psi_1$.
		At the cochain level this is $\diff \psi_2 = \Theta(\psi_1)$, for some
		$\Theta(\psi_1)\in \Cochain^2(G,M)$ with $\diff\Theta(\psi_1)=0$.
		Hence $\Theta(\psi_1)$ defines a class $[\Theta(\psi_1)]\in \HH^2(G,M)$.
		A solution $\psi_2$ exists if and only if $\Theta(\psi_1)$ is a coboundary,
		i.e.\ $[\Theta(\psi_1)]=0$.
	\end{proof}
	
	\begin{remark}[The obstruction as a Maurer--Cartan square]\label{rem:MC}
		The quadratic obstruction $\Theta(\psi_1)$ is the cohomological avatar of the self-interaction
		of the first-order term. Equip the cochain spaces $\Cochain^\bullet(G,M)$ with the cup-type product
		inherited from composition in $M=\End(E)$ and the resulting graded bracket $[\,\cdot,\cdot\,]$;
		then $(\Cochain^\bullet(G,M),\diff,[\,\cdot,\cdot\,])$ is a differential graded Lie algebra, and a
		direct computation gives $\Theta(\psi_1)=\tfrac12[f,f]$ for the classifying $1$-cocycle $f$. The
		identity $\diff\Theta=0$ is then the graded Jacobi relation, and the integrability of $f$ to a full
		formal deformation is precisely the Maurer--Cartan equation
		$\diff\xi+\tfrac12[\xi,\xi]=0$ for $\xi=tf+t^2\psi_2+\cdots$. This is the exact analogue, in the
		twin setting, of the role played by the Hochschild DGLA in associative-algebra deformation theory.
	\end{remark}

	\subsection{Rigidity}
	
	\begin{definition}[Infinitesimal rigidity]
		The twin system $\TwinSpace$ is \emph{infinitesimally rigid} if every infinitesimal
		$G$-equivariant deformation is equivalent to the trivial one.
	\end{definition}
	
	\begin{corollary}\label{cor:rigidity-HH1}
		$\TwinSpace$ is infinitesimally rigid if and only if $\HH^1(G,M)=0$.
	\end{corollary}
	
	\begin{proof}
		Immediate from \cref{thm:HH1-deformations}.
	\end{proof}
	
	%===========================================================
	\section{Deformation--Obstruction Theory for Twin Spaces}
	\label{sec:def-obs}
	
	We now formulate the full deformation--obstruction framework.
	
	\begin{definition}[Formal $G$-equivariant deformation]
		A \emph{formal $G$-equivariant deformation} of $\alpha$ is a formal power series
		\[
		\alpha_t(x,y) = \sum_{k\ge 0} t^k \alpha_k(x,y), \qquad \alpha_0=\alpha,
		\]
		such that for all $k\ge 0$,
		\[
		g\cdot\alpha_t(x,y) - \alpha_t(g\cdot x,g\cdot y) = O(t^{k+1}).
		\]
	\end{definition}
	
	Write $\alpha_t=\alpha + t\psi_1 + t^2\psi_2 + \cdots$. At each order we obtain
	\[
	\diff \psi_k = \Theta_k(\psi_1,\dots,\psi_{k-1}),
	\]
	where $\Theta_k$ is a universal expression, quadratic in the lower-order terms $\psi_1,\dots,\psi_{k-1}$ in the Gerstenhaber sense; in particular $\Theta_2$ depends only on $\psi_1$, quadratically.
	
	\begin{theorem}[Deformation--obstruction correspondence]\label{thm:def-obs}
		Let $\TwinSpace$ be a twin system, with $M=\End(E)$.
		\begin{enumerate}[label=(\roman*)]
			\item Equivalence classes of infinitesimal $G$-equivariant deformations
			are in bijection with $\HH^1(G,M)$.
			\item For each infinitesimal class $[f]\in \HH^1(G,M)$, the obstruction to extending it
			to a second-order deformation is a cohomology class in $\HH^2(G,M)$.
			\item A complete formal deformation (to all orders) exists if and only if all the
			obstruction classes vanish in $\HH^2(G,M)$ at each stage.
		\end{enumerate}
	\end{theorem}
	
	\begin{proof}
		Item~(i) is \cref{thm:HH1-deformations}; item~(ii) is \cref{thm:HH2-obstruction}.
		For~(iii), the higher-order extension steps give equations
		$\diff \psi_k = \Theta_k(\psi_1,\dots,\psi_{k-1})$ with $\diff\Theta_k=0$.
		At each stage the obstruction is a class in $\HH^2(G,M)$, and the existence of a complete formal solution
		is equivalent to the vanishing of all of these classes.
	\end{proof}
	
	\begin{definition}[Formal rigidity]
		The twin system $\TwinSpace$ is \emph{formally rigid} if any formal $G$-equivariant deformation
		is equivalent to the trivial one.
	\end{definition}
	
	\begin{theorem}[Cohomological rigidity]\label{thm:cohom-rigidity}
		If $\HH^1(G,M)=0$ and $\HH^2(G,M)=0$, then $\TwinSpace$ is formally rigid.
	\end{theorem}
	
	\begin{proof}
		If $\HH^1(G,M)=0$, every infinitesimal deformation is trivial.
		If in addition $\HH^2(G,M)=0$, there are no non-trivial obstruction classes and spectral
		sequences of obstructions collapse trivially. An induction on the order in $t$ shows that any
		putative deformation is equivalent to the trivial one.
	\end{proof}
	
	%===========================================================
	\section{Explicit Cohomology Computations}
	\label{sec:explicit}
	
	We now compute $\HH^\bullet(G,M)$ in several important cases.
	
	\subsection{Finite cyclic groups}
	
	Let $G=C_n=\langle g\mid g^n=e\rangle$ act semisimply on a finite-dimensional $k$-vector space $E$.
	The action extends to $M=\End(E)$.
	
	The bar resolution for $C_n$ is 2-periodic and yields standard formulas in terms of
	the operators
	\[
	(T-1) : M\to M,\quad A\mapsto g\cdot A - A,
	\]
	and the norm
	\[
	N : M\to M,\quad A\mapsto A + g\cdot A + \cdots + g^{n-1}\cdot A.
	\]
	
	Under mild genericity/semisimplicity assumptions (no non-trivial invariants except $M^G$), one has:
	
	\begin{theorem}\label{thm:cyclic}
		Let $G=C_n$ act semisimply on $E$ and $M=\End(E)$ as above.
		Then for all $k\ge 0$,
		\[
		\HH^{2k+1}(G,M) = 0, \qquad
		\HH^{2k}(G,M) \cong M^G.
		\]
		In particular, $\HH^1(G,M)=0$ and the corresponding twin system is infinitesimally rigid.
	\end{theorem}
	
	\begin{proof}
		From the 2-periodic resolution, one has
		\[
		H^{2k}(C_n,M) \cong \frac{\Ker(T-1)}{\Img(N)},\qquad
		H^{2k+1}(C_n,M) \cong \frac{\Ker(N)}{\Img(T-1)}.
		\]
		Semisimplicity implies the decomposition $M\simeq M^G\oplus M'$ with no invariants in $M'$.
		One checks that $\Ker(N)=\Img(T-1)$ and $\Ker(T-1)=M^G\oplus\Img(N)$,
		so odd cohomology vanishes and even cohomology identifies with $M^G$.
		Since our deformation cohomology $\HH^\bullet(G,M)$ coincides with this group cohomology,
		the result follows.
	\end{proof}
	
	\begin{example}[A rigid involution]\label{ex:C2-swap}
		Let $\operatorname{char}k\neq 2$, $E=k^2$, and $G=C_2=\{e,\sigma\}$ with $\sigma$ the coordinate
		swap $(x_1,x_2)\mapsto(x_2,x_1)$. On $M=\End(E)\cong M_2(k)$ the group acts by conjugation by the
		swap matrix $S=\bigl(\begin{smallmatrix}0&1\\1&0\end{smallmatrix}\bigr)$, namely $\sigma\cdot A=SAS^{-1}$.
		The invariants $M^G=\{A:SAS^{-1}=A\}$ are the matrices commuting with $S$, i.e.\ the $2$-dimensional
		centraliser $\{aI+bS:a,b\in k\}$. Since $\operatorname{char}k\neq 2$ the action is semisimple, so
		\cref{thm:cyclic} applies and gives
		\[
			\HH^{2k}(C_2,M)\cong M^G\;(\dim 2),\qquad \HH^{2k+1}(C_2,M)=0.
		\]
		In particular $\HH^1=0$: the associated twin system is infinitesimally rigid. The two-dimensional
		even cohomology is detected by the invariant endomorphisms $I$ and $S$ themselves.
	\end{example}

	\subsection{Finite abelian groups}
	
	Let $G$ be finite abelian and act semisimply on $E$ and $M=\End(E)$.
	We have a decomposition
	\[
	G \cong \Z/n_1\Z \times \cdots \times \Z/n_r\Z.
	\]
	Using Künneth, one obtains:
	
	\begin{theorem}\label{thm:abelian}
		Let $G$ be a finite abelian group acting semisimply on $E$.
		Then for all $k\ge 0$,
		\[
		\HH^k(G,M) \cong \bigwedge\nolimits^k \Hom(G,k)\,\otimes\, M^G.
		\]
		In particular,
		\[
		\HH^1(G,M)\cong \Hom(G,k)\otimes M^G.
		\]
	\end{theorem}
	
	\begin{proof}
		The cohomology $H^\bullet(G,k)$ is known to be an exterior algebra on $\Hom(G,k)$,
		and semisimplicity together with the universal coefficient theorem yields the stated
		tensor product with $M^G$.
	\end{proof}
	
	\subsection{Compact Lie groups}
	
	Let $G$ be a compact connected Lie group acting smoothly on a smooth manifold $E$.
	Consider $M=C^\infty(E)$ with the natural action $g\cdot f(x)=f(g^{-1}\cdot x)$.
	
	\begin{theorem}\label{thm:compact-vanish}
		If $G$ is compact connected and acts smoothly on $E$, then for all $k\ge 1$,
		\[
		\HH^k(G,C^\infty(E)) = 0.
		\]
	\end{theorem}
	
	\begin{proof}[Proof sketch]
		One uses Haar integration to define an averaging operator on cochains:
		\[
		(Af)(g_1,\dots,g_n)(x)
		= \int_G f(hg_1h^{-1},\dots,hg_nh^{-1})(h^{-1}\cdot x)\,d\mu(h),
		\]
		which projects onto conjugation-invariant cochains.
		For compact connected groups, such invariant cochains in the smooth setting
		are coboundaries in positive degrees, by the vanishing of suitable continuous group cohomology
		and the connectedness of $G$.
		Thus $\HH^k(G,C^\infty(E))$ vanishes for $k\ge 1$.
	\end{proof}
	
	\subsection{Matrix algebras and conjugation}
	
	Let $E=M_n(k)$ and $\alpha(X,Y)=XY$ be matrix multiplication. Let $G=GL_n(k)$ act by conjugation
	\[
	g\cdot X = gXg^{-1}.
	\]
	Then $M=\End(E)$ carries a highly semisimple representation.
	
	\begin{theorem}\label{thm:matrix-rigid}
		In the above setting,
		\[
		\HH^1(G,M) = 0, \qquad \HH^2(G,M) = 0.
		\]
		Hence the matrix twin system is formally rigid.
	\end{theorem}
	
	\begin{proof}[Proof sketch]
		Over a field of characteristic zero the group $GL_n$ is \emph{linearly reductive}: every rational
		representation splits as a direct sum of irreducibles, and the invariants functor $W\mapsto W^{GL_n}$
		is exact (it is split by the Reynolds operator). An exact invariants functor has no higher derived
		functors, so the rational cohomology of $GL_n$ vanishes in positive degree:
		$H^i(GL_n,W)=0$ for every rational module $W$ and every $i\ge 1$. Now $M_n(k)\cong V\otimes V^\ast$
		with $V=k^n$ under the conjugation action, so $M=\End(M_n(k))$ is again a rational $GL_n$-module;
		taking rational (polynomial) cochains---the natural choice for an algebraic group action---the
		vanishing applies to $W=M$ and yields $H^1=H^2=0$, hence
		$\HH^1$ and $\HH^2$ in our setting.
	\end{proof}
	
	%===========================================================
	\section{Spectral Sequences for Twin Deformation Cohomology}
	\label{sec:spectral}
	
	Let $1\to N\to G\to Q\to 1$ be a short exact sequence of groups.
	The action of $G$ on $M$ restricts to $N$ and descends to $Q$.
	
	\begin{theorem}[Hochschild--Serre spectral sequence]\label{thm:HS}
		There is a spectral sequence
		\[
		E_2^{p,q} = \HH^p(Q,\HH^q(N,M)) \Longrightarrow \HH^{p+q}(G,M).
		\]
	\end{theorem}
	
	\begin{proof}[Proof sketch]
		Filter the bar resolution of $G$ by the number of arguments in $N$ and $Q$, obtaining
		a double complex whose total complex computes $\HH^\bullet(G,M)$. The associated spectral
		sequence has the indicated $E_2$-page, as in~\cite{HochschildSerre1953}.
		Convergence follows from standard homological algebra arguments.
	\end{proof}
	
	\begin{example}[Dihedral groups]
		For the dihedral group $D_{2n}$ with normal cyclic subgroup generated by the rotation,
		the spectral sequence
		\[
		E_2^{p,q} = \HH^p(\Z/2\Z, \HH^q(\Z/n\Z,M))
		\]
		allows explicit computations of $\HH^\bullet(D_{2n},M)$ using the results of
		\cref{thm:cyclic,thm:abelian}.
	\end{example}
	
	%===========================================================
	\section{Rigidity Criteria for Twin Systems}
	\label{sec:rigidity}
	
	We summarize the cohomological criteria for rigidity.
	
	\begin{theorem}[Rigidity via $\HH^1$ and $\HH^2$]\label{thm:rigid-criteria}
		Let $\TwinSpace$ be a twin system with $M=\End(E)$.
		\begin{enumerate}[label=(\roman*)]
			\item $\TwinSpace$ is infinitesimally rigid if and only if $\HH^1(G,M)=0$.
			\item If $\HH^1(G,M)=0$ and $\HH^2(G,M)=0$, then $\TwinSpace$ is formally rigid.
		\end{enumerate}
	\end{theorem}
	
	\begin{proof}
		Item~(i) is \cref{cor:rigidity-HH1}. Item~(ii) is \cref{thm:cohom-rigidity}.
	\end{proof}
	
	\begin{corollary}[Compact connected groups]
		In the setting of \cref{thm:compact-vanish}, the associated twin system is formally rigid.
	\end{corollary}
	
	\begin{corollary}[Matrix twin systems]
		In the setting of \cref{thm:matrix-rigid}, the matrix twin system is formally rigid.
	\end{corollary}
	
	%===========================================================
	\section{Examples and Deformation Spectra}
	\label{sec:examples}
	
	We illustrate the variety of deformation behaviors.
	
	\begin{example}[Rigid cyclic rotations]
		Let $E=\R^d$ and $G=C_n$ act by a generic orthogonal rotation with no non-trivial fixed
		subspace except $\{0\}$. Then by \cref{thm:cyclic}, $\HH^1(G,M)=0$ and the twin system is
		infinitesimally rigid.
	\end{example}
	
	\begin{example}[A flexible free-abelian example]
		Take $G=\Z=\langle g\rangle$ acting on $E=k^2$ by $g=\operatorname{diag}(1,\mu)$ with $\mu\neq 0,1$, so that the
		conjugation action on $M=\End(E)\cong M_2(k)$ is semisimple. Group cohomology of $\Z$ is computed by the
		length-one free resolution $0\to\Z[\Z]\xrightarrow{\,g-1\,}\Z[\Z]\to\Z\to 0$, which gives
		$\HH^0=\ker(g-1)=M^G$, $\HH^1=\operatorname{coker}(g-1)=M/(g-1)M$, and $\HH^{\ge 2}=0$. The invariants
		$M^G$ are the matrices commuting with $\operatorname{diag}(1,\mu)$, namely the diagonal matrices, of
		dimension $2$; by semisimplicity $M=M^G\oplus(g-1)M$, so $\operatorname{coker}(g-1)\cong M^G$ and
		hence $\HH^1\cong M^G$ is also $2$-dimensional.
		Thus the twin system is genuinely \emph{flexible}: it carries a $2$-parameter family of pairwise
		inequivalent infinitesimal deformations, in agreement with the exterior-algebra computation underlying
		\cref{thm:abelian} applied to the free abelian group $\Z$ (for which $\Hom(\Z,k)=k$). A \emph{finite}
		abelian group over a characteristic-zero field would instead have $\Hom(G,k)=0$ and hence $\HH^1=0$;
		the flexibility here comes precisely from the infinite, free direction of $\Z$.
	\end{example}
	
	\begin{example}[Dihedral twin systems]
		For $D_{2n}$ acting on vertices or edges of a regular $n$-gon, the spectral sequence
		of \cref{thm:HS} combined with \cref{thm:cyclic} describes $\HH^\bullet(D_{2n},M)$,
		typically leading to non-trivial $\HH^1$ but constrained higher-degree cohomology.
	\end{example}
	
	\begin{example}[Compact actions on spheres]
		Let $G=SO(n)$ act on $E=S^{n-1}$. In many natural smooth settings, 
		\cref{thm:compact-vanish} implies vanishing of positive cohomology degrees, and hence
		rigidity of rotationally symmetric twin operations on the sphere.
	\end{example}
	
	%===========================================================
	\section{Conclusion and Outlook}
	\label{sec:conclusion}
	
	In this article we have developed a full deformation cohomology theory for \emph{twin spaces},
	i.e.\ triples $\TwinSpace$ equipped with the orbit of twin operations
	\[
	\alpha_g(x,y) = g^{-1}\cdot \alpha(g\cdot x,g\cdot y), \qquad g\in G,
	\]
	and the identification $\TwinOrbit \cong \LinearSubgroup\backslash G$ (right cosets) with the linear subgroup $\LinearSubgroup$.
	The key idea is that the relevant object is not the operation $\alpha$ alone,
	but its entire orbit of operations generated by the group action.
	This orbit structure is the conceptual basis for attaching a Hochschild-type cohomology complex
	$\Cochain^\bullet(G,\End(E))$.
	
	We have:
	\begin{itemize}
		\item constructed the complex $\Cochain^\bullet(G,\End(E))$ and defined $\HH^\bullet(G,\End(E))$;
		\item identified $\HH^0(G,\End(E))$ with invariant endomorphisms;
		\item shown that $\HH^1(G,\End(E))$ classifies infinitesimal $G$-equivariant deformations of the twin space;
		\item identified obstruction classes in $\HH^2(G,\End(E))$;
		\item established cohomological criteria for infinitesimal and formal rigidity;
		\item computed $\HH^\bullet(G,\End(E))$ in several families (finite cyclic and abelian groups, dihedral groups, compact Lie groups, matrix algebras);
		\item constructed a Hochschild--Serre spectral sequence for twin deformation cohomology.
	\end{itemize}
	
	Conceptually, twin spaces provide a bridge between symmetry, geometry and deformation:
	the group action organizes operations into a homogeneous orbit $\LinearSubgroup\backslash G$,
	and the way this space deforms is precisely encoded by the cohomology $\HH^\bullet(G,\End(E))$.
	In other words, the twin formalism is not a cosmetic reformulation; it is the structural
	reason why a rich deformation--obstruction theory exists in this setting.
	
	Possible directions for future work include:
	\begin{enumerate}[label=(\alph*)]
		\item \textbf{Categorical and operadic refinements.}
		In companion papers we promote twin spaces to objects of a 2-category endowed with monoidal
		and operadic structure. Understanding how $\HH^\bullet$ interacts with these higher structures
		is a natural next step.
		\item \textbf{Quantum and Hopf-theoretic applications.}
		Twin spaces can be used to reinterpret certain quantum deformations of algebras as deformations
		of transport rules on twin operations, suggesting a new viewpoint on quantum groups and Hopf algebras.
		\item \textbf{Derived and homotopical generalizations.}
		Replacing sets and vector spaces by chain complexes or higher stacks should lead to derived versions
		of twin spaces, with deformation controlled by higher or $\infty$-categorical analogues of
		$\HH^\bullet(G,\End(E))$.
		\item \textbf{Geometric and analytic implementations.}
		For actions on manifolds or analytic spaces, twin spaces offer a framework in which geometric constraints
		(compactness, regularity, curvature) translate directly into cohomological constraints on the admissible deformations.
	\end{enumerate}
	
	Altogether, twin spaces emerge as a robust and flexible foundation for encoding how operations transform
	under symmetry, with a deformation theory rich enough to connect to classical and modern themes in algebra,
	geometry, and mathematical physics.
	
	%===========================================================
	\begin{thebibliography}{99}
		
		\bibitem{Gerstenhaber1964}
		M.~Gerstenhaber,
		\newblock On the deformation of rings and algebras,
		\newblock {\em Ann. of Math. (2)} \textbf{79} (1964), 59--103.
		
		\bibitem{NijenhuisRichardson1967}
		A.~Nijenhuis and R.~W.~Richardson,
		\newblock Deformations of Lie algebra structures,
		\newblock {\em J. Math. Mech.} \textbf{17} (1967), 89--105.
		
		\bibitem{Drinfeld1986}
		V.~G.~Drinfeld,
		\newblock Quantum groups,
		\newblock in: {\em Proceedings of the International Congress of Mathematicians (Berkeley, 1986)}, 
		\newblock Amer. Math. Soc., 1987, pp.~798--820.
		
		\bibitem{Kontsevich2003}
		M.~Kontsevich,
		\newblock Deformation quantization of Poisson manifolds,
		\newblock {\em Lett. Math. Phys.} \textbf{66} (2003), 157--216.
		
		\bibitem{HochschildSerre1953}
		G.~Hochschild and J.-P.~Serre,
		\newblock Cohomology of group extensions,
		\newblock {\em Trans. Amer. Math. Soc.} \textbf{74} (1953), 110--134.
		
		\bibitem{Weibel1994}
		C.~A.~Weibel,
		\newblock {\em An Introduction to Homological Algebra},
		\newblock Cambridge Studies in Advanced Mathematics, 38, Cambridge Univ. Press, 1994.
		
	\end{thebibliography}
	
\end{document}
